\section{Selección del caso de negocio}

La propuesta seleccionada como caso de negocio para el desarrollo del proyecto final es el caso de TransMilenio S.A., únicamente el proceso de Gestión y despacho de flota desde el Centro de Control Integrado (CCI).

Entre las razones para escoger este caso, está el hecho de que el proceso depende directamente de información operacional que requiere ser recolectada, procesada y utilizada en tiempo real para apoyar con la supervisión, control y regulación de la flota. Dentro del proceso se manejan como la posición y velocidad de los buses, la programación de servicios, rutas, horarios, frecuencias, novedades operacionales y decisiones de regulación tomadas desde el CCI.

Al mirar este caso desde la óptica de un arquitecto de seguridad, la gestión y despacho de flota desde el CCI permite analizar aspectos relacionados con la tríada CIA, control de acceso y trazabilidad de la información que soporta la operación. Una afectación sobre alguno de estos elementos podría impactar directamente sobre la capacidad que tiene el CCI para tomar decisiones de despacho y regulación de manera óptima. Cabe resaltar que el proceso involucra diferentes actores, sistemas de información y componentes tecnológicos, suficiente para analizar los riesgos de seguridad de la información de manera articulada con los objetivos estratégicos de la organización.
